% 编译：在仓库根目录创建 output/pdf，然后将下列命令执行两遍。
% xelatex -halt-on-error -output-directory=output/pdf assignments/assignment10/assignment10-sol.tex
% 中文字体使用 TeX Live / MacTeX 自带的 Fandol 字体。
\documentclass[11pt,a4paper,fontset=fandol]{ctexart}
\usepackage[margin=2.1cm]{geometry}
\usepackage{amsmath,amssymb,mathtools}
\usepackage{fancyhdr,needspace}
\usepackage[hidelinks]{hyperref}
\hypersetup{pdftitle={线性代数第十次作业},pdfsubject={MATH1205：线性无关、极大线性无关组与秩}}
\setlength{\parindent}{0pt}
\setlength{\parskip}{0.25em}
\setlength{\headheight}{15pt}
\renewcommand{\arraystretch}{1.12}
\linespread{1.1}
\pagestyle{fancy}
\fancyhf{}
\fancyhead[L]{MATH1205\quad 线性代数}
\fancyhead[R]{第十次作业}
\fancyfoot[C]{\thepage}
\renewcommand{\headrulewidth}{0.3pt}
\ctexset{section={format=\large\bfseries,beforeskip=1.6ex,afterskip=0.7ex}}
\newcommand{\problem}[1]{\Needspace{6\baselineskip}\section*{#1}}
\newcommand{\ans}[1]{\boxed{#1}}
\newcommand{\rank}{\operatorname{r}}

\begin{document}
\setlength{\abovedisplayskip}{7pt plus 2pt minus 2pt}
\setlength{\belowdisplayskip}{7pt plus 2pt minus 2pt}
\setlength{\abovedisplayshortskip}{4pt plus 1pt minus 1pt}
\setlength{\belowdisplayshortskip}{4pt plus 1pt minus 1pt}
\begin{center}
{\LARGE\bfseries 线性代数第十次作业}\par
\vspace{0.3em}
{\small MATH1205}
\end{center}
\vspace{-0.3em}
按原 PDF 的页面顺序编号为第 1 至第 7 题。以下均在实数域上讨论。

\problem{第 1 题：循环相邻向量之和的线性相关性}
设 $\alpha_1,\ldots,\alpha_s\ (s\ge3)$ 线性无关，记
\[
\beta_i=\alpha_i+\alpha_{i+1}\quad(1\le i<s),\qquad
\beta_s=\alpha_s+\alpha_1.
\]
若 $k_1\beta_1+\cdots+k_s\beta_s=0$，按 $\alpha_i$ 合并得
\[
(k_1+k_s)\alpha_1+\sum_{i=2}^{s}(k_{i-1}+k_i)\alpha_i=0.
\]
由原向量组线性无关，
\[
k_1+k_s=0,\qquad k_i=-k_{i-1}\ (2\le i\le s).
\]
故 $k_i=(-1)^{i-1}k_1$，代入首式得
$\bigl(1+(-1)^{s-1}\bigr)k_1=0$。

当 $s$ 为奇数时，$2k_1=0$，从而所有 $k_i=0$，新向量组线性无关。
当 $s$ 为偶数时，有非平凡关系
\[
\beta_1-\beta_2+\beta_3-\cdots+\beta_{s-1}-\beta_s=0,
\]
所以新向量组线性相关。因此
\[
\ans{s\text{ 为奇数时线性无关；}\quad s\text{ 为偶数时线性相关。}}
\]

\problem{第 2 题：证明变换后的三个向量线性无关}
已知 $\alpha,\beta,\gamma$ 线性无关。设
\[
x(\alpha-\beta)+y(\beta+\gamma)+z(\gamma-\alpha)=0.
\]
整理得 $(x-z)\alpha+(-x+y)\beta+(y+z)\gamma=0$。
由原向量组线性无关，
\[
x-z=0,\qquad -x+y=0,\qquad y+z=0.
\]
前两式给出 $z=x,\ y=x$，代入第三式得 $2x=0$，故 $x=y=z=0$，
因此三个新向量\textbf{线性无关}。

\problem{第 3 题：两个向量相关与对应分量成比例的充要条件}
设 $n\ge2$，$\alpha=(a_1,\ldots,a_n)^T,\ \beta=(b_1,\ldots,b_n)^T$。
这里“对应分量成比例”指存在同一个实数 $\lambda$，使
$a_i=\lambda b_i$ 对所有 $i$ 成立，或 $b_i=\lambda a_i$ 对所有 $i$ 成立；
此表述也包含零向量的情形。

\textbf{必要性。} 若 $\alpha,\beta$ 线性相关，则存在不全为零的 $u,v$，
使 $u\alpha+v\beta=0$。
若 $u\ne0$，则 $\alpha=-\frac vu\beta$；
若 $u=0$，则 $v\ne0$，于是 $\beta=0=0\alpha$。
两种情形均说明对应分量成比例。

\textbf{充分性。} 若 $\alpha=\lambda\beta$，则
$\alpha-\lambda\beta=0$，系数 $(1,-\lambda)$ 不全为零；
若 $\beta=\lambda\alpha$，则同理由 $\beta-\lambda\alpha=0$ 得到非平凡关系。
故两向量线性相关，充要条件得证。

\newpage
\problem{第 4 题：能表示基本向量组的 $n$ 个向量线性无关}
设 $\varepsilon_1,\ldots,\varepsilon_n$ 是 $n$ 维基本向量组，且
\[
\varepsilon_j=\sum_{i=1}^n c_{ij}\alpha_i\qquad(j=1,\ldots,n).
\]
令 $A=(\alpha_1,\ldots,\alpha_n)$，$D=(c_{ij})_{n\times n}$。
将上述等式按列排列，得到
\[
AD=(\varepsilon_1,\ldots,\varepsilon_n)=I_n.
\]
两边取行列式，得
\[
\det A\cdot\det D=\det I_n=1.
\]
因此 $\det A\ne0$，矩阵 $A$ 的列向量线性无关，即
\[
\ans{\alpha_1,\ldots,\alpha_n\text{ 线性无关。}}
\]

\problem{第 5 题：四个四维向量的极大线性无关组与秩}
题中
\[
\alpha_1=\begin{pmatrix}3\\-5\\2\\1\end{pmatrix},\quad
\alpha_2=\begin{pmatrix}1\\1\\0\\-5\end{pmatrix},\quad
\alpha_3=\begin{pmatrix}-1\\3\\1\\3\end{pmatrix},\quad
\alpha_4=\begin{pmatrix}2\\-4\\-1\\-3\end{pmatrix}.
\]
以它们为列组成矩阵
\[
A=\begin{pmatrix}
3&1&-1&2\\
-5&1&3&-4\\
2&0&1&-1\\
1&-5&3&-3
\end{pmatrix}.
\]
对行列式依次作列变换
$C_1\leftarrow C_1-3C_2$、$C_3\leftarrow C_3+C_2$、
$C_4\leftarrow C_4-2C_2$，其值不变。沿第一行展开，得到
\[
\begin{aligned}
\det A
&=\begin{vmatrix}
0&1&0&0\\-8&1&4&-6\\2&0&1&-1\\16&-5&-2&7
\end{vmatrix}
=-\begin{vmatrix}-8&4&-6\\2&1&-1\\16&-2&7\end{vmatrix}\\[2pt]
&=-\begin{vmatrix}0&8&-10\\2&1&-1\\0&-10&15\end{vmatrix}
=2\begin{vmatrix}8&-10\\-10&15\end{vmatrix}
=2(120-100)=40\ne0.
\end{aligned}
\]
其中三阶行列式用了
$R_1\leftarrow R_1+4R_2$、$R_3\leftarrow R_3-8R_2$。
故四个向量线性无关，整组就是极大线性无关组，且
\[
\ans{\text{极大线性无关组为 }\{\alpha_1,\alpha_2,\alpha_3,\alpha_4\},
\qquad \rank(A)=4.}
\]

\newpage
\problem{第 6 题：求秩、极大线性无关组及其余向量的表示}
题中
\[
\alpha_1=\begin{pmatrix}2\\4\\2\end{pmatrix},\quad
\alpha_2=\begin{pmatrix}1\\1\\0\end{pmatrix},\quad
\alpha_3=\begin{pmatrix}2\\3\\1\end{pmatrix},\quad
\alpha_4=\begin{pmatrix}3\\5\\2\end{pmatrix}.
\]
将向量按列排成矩阵，作初等行变换：
\[
\begin{aligned}
A&=\begin{pmatrix}2&1&2&3\\4&1&3&5\\2&0&1&2\end{pmatrix}
\xrightarrow[\ R_3\leftarrow R_3-R_1\ ]{R_2\leftarrow R_2-2R_1}
\begin{pmatrix}2&1&2&3\\0&-1&-1&-1\\0&-1&-1&-1\end{pmatrix}\\[3pt]
&\xrightarrow{R_3\leftarrow R_3-R_2}
\begin{pmatrix}2&1&2&3\\0&-1&-1&-1\\0&0&0&0\end{pmatrix}
\sim
\begin{pmatrix}1&0&\frac12&1\\0&1&1&1\\0&0&0&0\end{pmatrix}.
\end{aligned}
\]
最后一步依次使用 $R_2\leftarrow-R_2$、$R_1\leftarrow R_1-R_2$、
$R_1\leftarrow\frac12R_1$。非零行有两行，主元在第 $1,2$ 列，因此
\[
\ans{\rank(A)=2,\qquad \text{一个极大线性无关组为 }\{\alpha_1,\alpha_2\}.}
\]
初等行变换保持列向量之间的线性关系，由最后的矩阵读出
\[
\ans{\alpha_3=\frac12\alpha_1+\alpha_2,\qquad
\alpha_4=\alpha_1+\alpha_2.}
\]
代回验证：
\[
\frac12\begin{pmatrix}2\\4\\2\end{pmatrix}
+\begin{pmatrix}1\\1\\0\end{pmatrix}
=\begin{pmatrix}2\\3\\1\end{pmatrix},\qquad
\begin{pmatrix}2\\4\\2\end{pmatrix}
+\begin{pmatrix}1\\1\\0\end{pmatrix}
=\begin{pmatrix}3\\5\\2\end{pmatrix}.
\]

\problem{第 7 题：横向拼接矩阵的秩不等式}
设 $A,B$ 均为 $m\times n$ 矩阵，$C=(A,B)$。记
$\rank(A)=p,\ \rank(B)=q$。

\textbf{下界。} $A$ 和 $B$ 的列向量组都是 $C$ 的列向量组的子组，
故其中的线性无关组仍是 $C$ 中的线性无关组。因此
\[
\rank(A)\le\rank(C),\qquad \rank(B)\le\rank(C),
\quad\text{即}\quad \max\{\rank(A),\rank(B)\}\le\rank(C).
\]

\textbf{上界。} 分别从 $A,B$ 的列向量组中取极大线性无关组
$\xi_1,\ldots,\xi_p$ 和 $\eta_1,\ldots,\eta_q$。
$C$ 的每一列均可由这两组共 $p+q$ 个向量线性表示，
所以 $\rank(C)\le p+q$。当某个秩为零时，相应组取空组，论证仍成立。

综上，
\[
\ans{\max\{\rank(A),\rank(B)\}\le\rank(C)\le\rank(A)+\rank(B).}
\]
\end{document}
